\section{Conclusiones generales}
\label{sec:concl_generales}
% Extensión estimada: 1.5 páginas, cinco párrafos (ficha 7.1).
% Fuentes: sec:res_clasif_discusion, sec:res_gen_discusion, sec:disc_representacion, sec:disc_texto_imagen,
% sec:disc_concordancia, sec:disc_limites. Ajustes acordados con el autor en la reestructuración del Cap. 6: PI1
% «depende del uso», no «negativa»; la comparación de DermLIP con una cabeza es entre dos formas de aprender la
% representación y no se sometió a prueba; PI2 sin tono tajante, con la verificación parcial de la descripción.

El trabajo partió de dos preguntas: si los \glspl{vlm} clasifican las lesiones de piel tan bien como los clasificadores
entrenados solo con imágenes, y si los reportes que genera un \gls{vlm} afirman el diagnóstico correcto y describen la
lesión de la imagen. Las dos se respondieron sobre la prueba de \emph{HAM10000} con una partición disjunta por lesión,
con métricas automáticas y con el cómputo de una sola GPU, y sus respuestas valen dentro de esas condiciones.

La respuesta a la primera pregunta depende de cómo se use el \gls{vlm}. Como clasificador autónomo, ninguno de los
evaluados alcanza a los supervisados. Los generativos de propósito general quedan cerca del azar, porque su respuesta
depende sobre todo del \emph{prompt}, y DermLIP sin ajuste queda entre 0.15 y 0.18 de sensibilidad macro por debajo de
los mejores supervisados, con diferencias que las pruebas estadísticas respaldan
(Capítulo~\ref{cp:res_clasif}). Como codificador de un modelo especializado, en cambio, se acerca: con una cabeza
entrenada sobre la representación congelada de DermLIP, el F1 macro llega a 0.705, frente a 0.733 de Swin-Tiny, una
diferencia que no se sometió a prueba. Esa cercanía indica que la especialización en dermatología está en la
representación, pero no cumple la motivación de la pregunta, que era obtener el diagnóstico del mismo modelo que
escribe el reporte (Sección~\ref{sec:disc_representacion}). En el melanoma, la clase de mayor prioridad clínica, la
diferencia la marca la regla de decisión más que el modelo: el umbral fijado en validación lleva la detección de
Swin-Tiny de tres a cuatro de cada cinco casos.

La respuesta a la segunda pregunta es que el reporte afirma el diagnóstico correcto en la medida en que lo hace el
clasificador que lo ancla. Sin un diagnóstico, ningún generador afirma de manera confiable la clase de la imagen: el
que más la nombra lo hace enumerando dos clases por reporte, y en ninguno más de uno de cada veinte reportes nombra
únicamente la clase verdadera. Con la cascada, en cambio, el generador transcribe el diagnóstico que recibe, de modo que la consistencia
diagnóstica de Mini-VLM, 0.849, coincide con la exactitud de Swin-Tiny, 0.847, y con un diagnóstico equivocado el
reporte lo afirma en lugar de corregirlo (Capítulo~\ref{cp:res_gen}). El ruteo por concordancia eleva la consistencia
a 0.870 y, entre los diagnósticos que afirma, reduce la proporción de errores de 0.153 a 0.085
(Sección~\ref{sec:disc_concordancia}). En cuanto a la descripción, la imagen cambia lo que el reporte describe de cada
lesión, porque sin ella la cobertura de entidades clínicas de Mini-VLM cae a la mitad, pero la descripción también se
adapta al diagnóstico recibido, y que cada rasgo descrito esté en la imagen queda verificado solo en parte.

Las dos respuestas se encadenan. Como un \gls{vlm} generativo no decide la clase por sí mismo, el diagnóstico lo pone
un clasificador; y como el generador lo transcribe, la fidelidad diagnóstica del reporte queda fijada por ese
clasificador y no por el generador (Sección~\ref{sec:disc_texto_imagen}). En consecuencia, el valor del generador
está en la redacción y en la descripción de cada lesión, y el de la concordancia con el \gls{knn}, en señalar los casos
en que el diagnóstico es menos confiable. Lo que el sistema no resuelve son los errores que comparten los dos
predictores, que se afirman con la misma plantilla que un acierto y que incluyen casi uno de cada cuatro melanomas. Por
eso el \emph{pipeline} puede ordenar la atención del especialista, pero no reemplazarla.

Por último, el trabajo muestra que las decisiones de evaluación pesan tanto como la elección del modelo. La partición
por imagen sobreestima la sensibilidad macro en unas 0.07, una cifra del orden de la distancia con los mejores
resultados de la literatura, y la variación entre ejecuciones supera una décima en las clases minoritarias, mientras
que Swin-Tiny, ConvNeXt-Tiny y ViT-B/16 no se distinguen entre sí (Sección~\ref{sec:disc_limites}). Esa observación,
junto con el costo medido, hace compatible la restricción de una GPU de consumo con la fidelidad diagnóstica del
\emph{pipeline}: Mini-VLM ocupa como máximo 8.6~GB y escribe un reporte en unos 16 segundos, y la fidelidad depende de
un clasificador pequeño.
